\section{Computations}\label{sec:verify}

No proof uses a computed number. The computations check the exact identities and test
the theorems at moderate $N$. Except where we say otherwise they use the zeta law
$f(n)=n^{-\alpha-1}/\zeta(\alpha+1)$ for both signs, so $b=1$, and the fair start. The
exact law comes from~\eqref{eq:run-sum}. The convolution powers $u^\pm_m$ are computed by
direct convolution, and every entry of $A$, $B$ and $F^\plus_N$ is a sum of nonnegative
terms, so the relative rounding error stays small. The total mass is $1$ to $12$ digits in
every case. The identities of Section~\ref{sec:identities} and~\eqref{eq:run-sum} agree
with a sum over all sign strings for $N\le12$, for two different laws of the two signs and
a general initial run, to $2\cdot10^{-16}$.

Table~\ref{tab:ratios} shows the ratio $NP_N(k)/L_{\alpha,1}(k/N)$ at $k=N/2$ and $k=N/20$,
and the number of sign changes of $k\mapsto P_N(k+1)-P_N(k)$. At $\alpha=1/2$ the central
ratio is $1.0149$, $1.0105$ and $1.0075$ for $N=1000$, $2000$ and $4000$, and the
distance to $1$ shrinks by a factor close to $\sqrt2$ when $N$ doubles. The convergence is
slower away from $\alpha=1/2$ and much slower as $\alpha\to1$. At $\alpha=0.9$ the central
ratio is $0.54$ at $N=4000$. The sign changes sit at $k=N/2$ for $\alpha\le0.58$. At
$\alpha=0.7$ and $N=1000$ they sit at $k=104$, $500$ and $896$, near the minima of
$L_{0.7,1}$. At $\alpha=0.9$ and $N=4000$ they sit at $k=1$, $142$, $2000$, $3858$ and
$3999$. The two outer ones come from the ends, as in Example~\ref{ex:counts}. Indeed
$P_N(N)/T(N)=\frac12$ exactly, while $P_N(N-1)/T(N)$ is $0.5754$, $0.5737$ and $0.5727$ for
$N=1000$, $2000$ and $4000$, tending to $f(1)=0.5715$. On both sides of the threshold of
Theorem~\ref{thm:shape} the count is $1$ at $\alpha=0.58$ and $3$ at $\alpha=0.61$, for all
three $N$. Every count in the table equals the lower bound of Theorem~\ref{thm:signs}.

\begin{table}[t]
\centering
\begin{tabular}{@{}cccccc@{}}
\toprule
$\alpha$ & $N$ & ratio at $k=N/2$ & ratio at $k=N/20$ & sign changes \\
\midrule
$0.3$ & $1000$, $2000$, $4000$ & $1.0852$, $1.0687$, $1.0555$ & $1.1023$, $1.0845$, $1.0694$ & $1$, $1$, $1$\\
$0.5$ & $1000$, $2000$, $4000$ & $1.0149$, $1.0105$, $1.0075$ & $1.0045$, $1.0036$, $1.0028$ & $1$, $1$, $1$\\
$0.58$ & $1000$, $2000$, $4000$ & $0.9846$, $0.9863$, $0.9883$ & $1.0142$, $1.0114$, $1.0090$ & $1$, $1$, $1$\\
$0.61$ & $1000$, $2000$, $4000$ & $0.9696$, $0.9739$, $0.9781$ & $1.0282$, $1.0230$, $1.0185$ & $3$, $3$, $3$\\
$0.7$ & $1000$, $2000$, $4000$ & $0.9023$, $0.9159$, $0.9285$ & $1.1202$, $1.1004$, $1.0833$ & $3$, $3$, $3$\\
$0.9$ & $1000$, $2000$, $4000$ & $0.4903$, $0.5168$, $0.5433$ & $2.2892$, $2.1277$, $1.9905$ & $5$, $5$, $5$\\
\bottomrule
\end{tabular}
\caption{The ratio $NP_N(k)/L_{\alpha,1}(k/N)$ and the number of sign changes of the first
differences of $P_N$, for the zeta law and the fair start.}
\label{tab:ratios}
\end{table}

At the end, for $\alpha=1/2$, the ratio $P_N(N-j)/(T(N)u(j))$ is given in
Table~\ref{tab:ends} for $1\le j\le5$. It lies in $[1.0009,1.0040]$ at $N=1000$ and in
$[1.0004,1.0020]$ at $N=2000$, so the error halves when $N$ doubles.

\begin{table}[t]
\centering
\begin{tabular}{@{}lccccc@{}}
\toprule
& $j=1$ & $j=2$ & $j=3$ & $j=4$ & $j=5$\\
\midrule
$N=1000$ & $1.0009$ & $1.0018$ & $1.0026$ & $1.0033$ & $1.0040$\\
$N=2000$ & $1.0004$ & $1.0009$ & $1.0013$ & $1.0017$ & $1.0020$\\
\bottomrule
\end{tabular}
\caption{The ratio $P_N(N-j)/(T(N)u(j))$ of Theorem~\ref{thm:ends}(a) for the zeta law with
$\alpha=1/2$ and the fair start.}
\label{tab:ends}
\end{table} The ratio of $P_N(N-j)$ to $T(N)u(j)\Phi_1(j/N)$ over $1\le j\le0.8N$ lies
in $[1.0004,1.0148]$ at $N=1000$ and in $[1.0002,1.0105]$ at $N=2000$, as
Theorem~\ref{thm:ends}(b) predicts.

For unequal laws we take the zeta law for the plus runs and
$f_\minus(n)\propto(n+3)^{-\alpha-1}$ for the minus runs, so $b=\zeta(\alpha+1)/\zeta(\alpha+1,4)$,
with a fresh plus start. At $\alpha=1/2$, where $b=2.4498$, the ratios to $L_{\alpha,b}$ at
$x=0.1,0.3,\dots,0.9$ lie in $[1.0121,1.0397]$ at $N=1000$ and in $[1.0069,1.0199]$ at
$N=4000$. At $\alpha=0.7$, where $b=3.4700$, they lie in $[0.8460,1.2994]$ at $N=1000$ and in
$[0.8855,1.1920]$ at $N=4000$. In both cases the number of sign changes is $2$, at $k=1$ and at
the interior minimum, as in Example~\ref{ex:counts}.

For the Lamperti density, the root of $\cos(\pi\alpha/2)=\alpha$ is
$0.594611644056836$, and on a grid of $2\cdot10^5$ points $L_{\alpha,1}$ has $1$ critical
point for $\alpha\le0.5946$ and $3$ for $\alpha\ge0.5947$, at the values we tried. The
transform~\eqref{eq:stieltjes} agrees with direct quadrature in $30$-digit arithmetic, and
at $\alpha=1/2$ the double integral~\eqref{eq:Lambda} agrees with $L_{1/2,1}$ to a relative
error below $10^{-13}$ at $x=0.2$, $0.5$ and $0.7$.

For the initial run we take $\alpha=1/2$, a fair sign and $\nu_\pm(d)=d^{-1.2}/(2\zeta(1.2))$.
Then $n\nu_\pm(n)\to0$, and Theorem~\ref{thm:initial}(c) gives the local law. The central
ratio is $0.815$, $0.836$ and $0.855$ for $N=500$, $1000$ and $2000$. The gap to $1$ decays
like $N^{-0.2}$, the order of $\Prob(D\ge N)$, which enters through the first run with
$\varepsilon N\le D<N$ in the proof of Theorem~\ref{thm:initial}(a). For $D$ uniform on
$\{1,\dots,N/2\}$ with a fair sign, $NP_N(k)$ at $k/N=0.3$ and $0.5$ is $0.898$ and $1.325$
at $N=500$ and $0.884$ and $1.299$ at $N=2000$. The formula of Conjecture~\ref{rem:aging}
gives $0.8694$ and $1.2732$, against the values $0.6946$ and $0.6366$ of $L_{1/2,1}$ for a
fresh start.

The scripts and their outputs, each headed by the hash of its script, are in the folder
\nolinkurl{kagey-problems/problem131/paperG} of \url{https://github.com/apemm/Kagey-Problems}.
The command \texttt{python -I verify\_all.py} runs all checks in about a minute, and
\nolinkurl{README.md} lists what each script tests.
